\chapter{Black-Scholes Model}



\section{Closed-form solution for European options}

\section{Binomial Tree Method}
[Need to make more cohesive]

To price the put option with the tree method, we will use the binomial discrete-time model. This is an abstraction of the market dynamics, where the stock price can only move to one of two possible new prices in a time step. This is a simple model which we can make recombinant so that the algorithm performs with quadratic time complexity, which will serve as a benchmark for our numerical results for other methods.

To formalise the tree construction, let’s fix a time horizon $T > 0$. For some depth $n \ge 1$, divide the interval $[0, T]$ into $n$ even steps each of length $\Delta t = T/n$. Our asset $S_t$ will follow a multiplicative binomial evolution:
\begin{equation}
    S_{i+1} = \begin{cases}
    S_i u & \text{with probability }p \\S_i d & \text{with probability }1-p
    \end{cases}\qquad S_0 > 0
\end{equation}
where $u$ and $d$ are the model parameters determining the two factors of the stock price change per time step, where we choose $u > 1$ and $d \in (0,1)$.

The financial market will also consist of a risk-free asset $B_t$ satisfying
\begin{equation}
    B_{i+1} = B_i e^{r\Delta t} \qquad B_0 = 1
\end{equation}
where $r$ is the continuously-compounded risk-free rate, typically considered interest. For an arbitrage free market, we should additionally ensure that
\begin{equation}
u > e^{r\Delta t} > d
\end{equation}
for sufficiently small $\Delta t$, else we can make a riskless profit by shorting the stock and investing in the risk-free asset, or vice versa.

To ensure the absence of arbitrage, we require the discounted asset price to have zero drift (admitting a martingale measure). More precisely, let the discounted price be
\begin{equation}
\tilde S_i = e^{-ri\Delta t} S_i
\end{equation}
We require some probability measure $\mathbb Q$, equivalent to the physical measure $\mathbb P$ (have the same non-zero support) such that
\begin{equation}
\mathbb E_\mathbb Q [\tilde S_{i+1} | \mathcal F_i] = \tilde S_i
\end{equation}
To calculate the risk-neutral probability, we can apply this to our model:
\begin{equation}
e^{-r(i+1)\Delta t}\big[pS_i u + (1-p)S_i d \big] = e^{-ri\Delta t}S_i
\end{equation}

Collect the $p$ terms and cancel out the discount factor to get
\begin{equation}
p(S_i u - S_i d) + S_i d = e^{r\Delta t} S_i
\end{equation}

Divide both sides by $S_i$ and rearrange to make $p$ the subject to obtain
\begin{equation}
p = \frac{e^{r \Delta t} -d}{u-d}
\end{equation}
which we will use in our binomial model.

We will adopt a backward induction approach to price the American put at time $0$, since we already know the price at maturity, which is simply the payoff $g: \mathbb R^+ \to \mathbb R$ of the put option:

\begin{equation}
g(S_T) = (K - S_T)^+
\end{equation}

As the discounted price process is a martingale under the risk-free measure $\mathbb Q$, the arbitrage-free value $v(i\Delta t)$ of the derivative at time $i \Delta t$ is necessarily

\begin{equation}
v_i = \underbrace{e^{-r(T-i\Delta t)}}_\text{discount} \,\mathbb E_\mathbb Q[g(S_T)\,|\,\mathcal F_i]
\end{equation}

We will utilise the Cox-Ross-Rubinstein (CRR) Model \cite{cox1979option} for our implementation, which chooses the parameters
\begin{equation}
u = e^{\sigma \sqrt{\Delta t}}\qquad d = e^{-\sigma \sqrt{\Delta t}}
\end{equation}
satisfying the arbitrage free condition, and normalising the local variance of returns to the volatility parameter $\sigma^2$. They also conveniently multiply to $1$.

To implement this, we need to first compute the price tree. This would be, for the price at step $i$ after $j$ up-moves (so clearly $i \ge j$)
\begin{equation}
S_{i, j} = S_0 u^j d^{i-j}
\end{equation}

As the payoff is known at maturity, we can simply set
\begin{equation}
v_n(j) = g(S_{n, j}) \qquad 0 \le j \le n
\end{equation}
and recursively set
\begin{equation}
v_i(j) = e^{-r\Delta t} \big(p^* v_{i+1}(j+1) + (1-p^*)v_{i+1}(j)\big)
\end{equation}
for $0 \le j \le i$. This would be it for European options, but since early exercise is allowed, we will compare this value with the payoff of the price at step $i$ with $j$ up-moves, which is just $f(S_{i, j})$. Thus the value is given by
\begin{equation}
v_i(j) = \max\bigg(e^{-r\Delta t} \big(p^* v_{i+1}(j+1) + (1-p^*)v_{i+1}(j)\big), \; g(S_{i, j})\bigg)
\end{equation}
The root value $v_0(0)$ would be the arbitrage-free option price.

\section{Physics-Informed Neural Networks}

\section{Sobolev Deep Neural Networks}